% Conservative Force Proof for the Casimir Interaction
%
% ASSERT_CONVENTION: natural_units=natural, metric_signature=mostly_minus,
%   fourier_convention=physics, regularization_scheme=zeta
% ASSERT_CONVENTION: casimir_force_sign=F < 0 for attractive
% ASSERT_CONVENTION: casimir_energy_sign=E < 0 for bound configuration
% ASSERT_CONVENTION: work_sign=W_net > 0 means extraction; no-go claim: W_net <= 0
% ASSERT_CONVENTION: [F/A] = length^{-4}, [E/A] = length^{-3}, [W/A] = length^{-3}
%
% Reference: H.B.G. Casimir, Proc. K. Ned. Akad. Wet. 51, 793 (1948)
%   E/A = -pi^2/(720 a^3), F/A = -pi^2/(240 a^4) for ideal parallel plates at T=0

\documentclass[11pt]{article}
\usepackage{amsmath,amssymb,amsthm}
\usepackage[margin=1in]{geometry}

\newtheorem{theorem}{Theorem}
\newtheorem{corollary}{Corollary}
\newtheorem{definition}{Definition}

\title{Conservative Nature of the Casimir Force\\and Vanishing Net Work in Quasi-Static Cycles}
\author{GPD Research Project --- Phase 03, Plan 01}
\date{March 2026}

\begin{document}
\maketitle

%% ============================================================
\section{Conventions and Definitions}
%% ============================================================

Throughout, we work in natural units $\hbar = c = k_B = 1$.
All quantities are per unit plate area $A$ unless otherwise noted.
The metric signature is $(+,-,-,-)$.

\begin{definition}[Sign conventions]
\begin{itemize}
  \item \textbf{Force:} $F < 0$ for attractive (plates pulled together).
  \item \textbf{Energy:} $E < 0$ for a bound configuration (energy lower than infinite separation).
  \item \textbf{Work:} $W > 0$ when the Casimir force does positive work on the plates
    (energy is transferred \emph{from} the vacuum field to the mechanical degree of freedom).
    $W_\mathrm{net} > 0$ would mean net energy extraction over a cycle.
    The no-go claim is $W_\mathrm{net} \le 0$.
\end{itemize}
\end{definition}

%% ============================================================
\section{Theorem: The Casimir Force is Conservative}
%% ============================================================

\begin{theorem}[Conservative Casimir force]
\label{thm:conservative}
For parallel plates with fixed dielectric response functions
$\varepsilon(\omega)$ and $\mu(\omega)$, the Casimir force
\begin{equation}
  F(a,T) = -\frac{\partial V(a,T)}{\partial a}
  \label{eq:force-potential}
\end{equation}
derives from a single-valued potential $V(a,T)$, where:
\begin{itemize}
  \item At $T = 0$: $V(a) \equiv E_\mathrm{vac}(a)$, the vacuum energy.
  \item At $T > 0$: $V(a,T) \equiv \mathcal{F}(a,T) = -T \ln Z(a,T)$,
    the Helmholtz free energy.
\end{itemize}
The force $F$ is therefore conservative under conditions A1--A4 (below),
and the work done over any closed path in configuration space vanishes:
\begin{equation}
  W_\mathrm{net} = \oint F(a)\, da = 0.
  \label{eq:cycle-vanish}
\end{equation}
\end{theorem}

%% ============================================================
\section{Conditions for Conservativeness (A1--A4)}
%% ============================================================

The force $F = -\partial V / \partial a$ is conservative if and only if
the potential $V$ is a single-valued function of the configuration.
This holds under the following conditions:

\begin{description}
  \item[A1 --- Fixed material response:]
    The dielectric function $\varepsilon(\omega)$ and magnetic
    permeability $\mu(\omega)$ are \emph{independent of plate separation}
    and remain unchanged during the cycle. This ensures $V$ depends on $a$
    (and $T$) only through the mode structure, not through a hidden
    dependence on material state.

  \item[A2 --- Quasi-static process:]
    The system is in thermal equilibrium (or the vacuum state, at $T=0$)
    at every intermediate step. The plate velocity satisfies $v \ll c$,
    so no dynamical Casimir photons are produced. The force at each $a$
    is the equilibrium force $F(a,T)$.

  \item[A3 --- Isothermal process:]
    The temperature $T$ is held fixed throughout the cycle, OR the
    material response $\varepsilon(\omega)$, $\mu(\omega)$ is
    $T$-independent so that temperature variations do not alter the
    force law. For an isothermal process, $V(a,T) = \mathcal{F}(a,T)$
    is the appropriate potential.

  \item[A4 --- Fixed geometry class:]
    The plates remain parallel throughout the cycle. No topology change
    (e.g., merging plates, creating/destroying a cavity) occurs.
    The configuration space is $a \in (0, \infty)$, a simply-connected
    one-dimensional manifold, on which any single-valued $V(a)$ is
    automatically conservative.
\end{description}

%% ============================================================
\section{Proof}
%% ============================================================

\subsection{Case 1: $T = 0$ (vacuum energy)}

At $T = 0$, the Casimir energy is given by the regularized vacuum energy
of the electromagnetic field between the plates.
For ideal perfectly conducting plates (Casimir 1948):
\begin{equation}
  \frac{E(a)}{A} = -\frac{\pi^2}{720\, a^3}.
  \label{eq:casimir-energy}
\end{equation}

% DIMENSION CHECK: [E/A] = [a^{-3}] = [length^{-3}]. CHECK.

The function $E(a)$ depends on $a$ only (no other dynamical variable
for fixed plate material at $T=0$).
It is manifestly single-valued on $(0,\infty)$.

The force per unit area is:
\begin{equation}
  \frac{F(a)}{A} = -\frac{d}{da}\frac{E(a)}{A}
  = -\frac{\pi^2}{720}\cdot(-3)\cdot a^{-4}
  = -\frac{\pi^2}{240\, a^4}.
  \label{eq:casimir-force}
\end{equation}

% SIGN CHECK: d/da(a^{-3}) = -3 a^{-4}. Minus sign from F = -dE/da
% gives F = -(-3)(-pi^2/720) a^{-4} = -(3 pi^2/720) a^{-4} = -pi^2/(240 a^4).
% F < 0: attractive. CHECK.
%
% DIMENSION CHECK: [F/A] = [a^{-4}] = [length^{-4}]. CHECK.

Since $V(a) = E(a)$ is single-valued and $F = -dV/da$, the force is
conservative. \qed

\subsection{Case 2: $T > 0$ (Helmholtz free energy)}

At finite temperature $T > 0$, the thermal equilibrium state of the
field-plus-plates system is described by the partition function $Z(a,T)$.
The Helmholtz free energy
\begin{equation}
  \mathcal{F}(a,T) = -T \ln Z(a,T)
  \label{eq:helmholtz}
\end{equation}
is a \emph{state function}: it depends on the thermodynamic state
$(a, T)$ only, not on the path by which the state was reached.

Under conditions A1--A4:
\begin{itemize}
  \item $Z(a,T)$ is a function of $(a,T)$ only (no hidden variables),
    because the Hamiltonian depends on $a$ through the boundary conditions
    and the material response functions are fixed (A1).
  \item $\mathcal{F}$ is single-valued on the configuration space.
  \item The force $F = -\partial \mathcal{F}/\partial a |_T$ is therefore
    the gradient of a state function at fixed $T$, hence conservative.
\end{itemize}

The Lifshitz formula gives the force explicitly:
\begin{equation}
  \frac{F(a,T)}{A} = -\frac{T}{\pi}\,
  \sideset{}{'}\sum_{l=0}^{\infty}\;
  \int_0^\infty dk_\perp\, k_\perp\, \kappa_0 \,
  \sum_{P=\mathrm{TE,TM}}
  \frac{r_P^2\, e^{-2\kappa_0 a}}{1 - r_P^2\, e^{-2\kappa_0 a}}
  \label{eq:lifshitz-force}
\end{equation}
where $\xi_l = 2\pi l T$ are the Matsubara frequencies,
$\kappa_0 = \sqrt{k_\perp^2 + \xi_l^2}$, and $r_P$ are the
Fresnel reflection coefficients at imaginary frequency.
This is manifestly a function of $(a, T)$ only (for fixed
$\varepsilon(\omega)$, $\mu(\omega)$), confirming single-valuedness of
the underlying potential $\mathcal{F}(a,T)$.

Under an isothermal quasi-static process (A2, A3), the relevant
thermodynamic potential is $\mathcal{F}(a,T)$, and the force is its
gradient. \qed

%% ============================================================
\section{Corollary: Vanishing Cycle Work}
%% ============================================================

\begin{corollary}[Zero net work in closed cycle]
\label{cor:vanishing-work}
For any closed quasi-static isothermal path $a_1 \to a_2 \to a_1$
satisfying conditions A1--A4:
\begin{equation}
  W_\mathrm{net} = \oint F(a)\, da
  = \int_{a_1}^{a_2} F(a)\, da + \int_{a_2}^{a_1} F(a)\, da = 0.
  \label{eq:wnet-zero}
\end{equation}
\end{corollary}

\begin{proof}
Since $F = -dV/da$ for a single-valued potential $V(a)$:
\[
  \int_{a_1}^{a_2} F(a)\,da
  = -\int_{a_1}^{a_2} \frac{dV}{da}\, da
  = -(V(a_2) - V(a_1))
  = V(a_1) - V(a_2).
\]
Similarly, $\int_{a_2}^{a_1} F(a)\, da = V(a_2) - V(a_1)$.
Adding: $W_\mathrm{net} = 0$. \qed
\end{proof}

This result is exact: it follows from the existence of a potential,
not from any approximation. It holds for ideal plates at $T=0$,
for realistic materials at finite $T$, and for \emph{any} cycle
geometry $a_\mathrm{min} \to a_\mathrm{max} \to a_\mathrm{min}$,
provided A1--A4 are satisfied.

%% ============================================================
\section{Analytical $T = 0$ Ideal-Plate Cycle}
%% ============================================================

\subsection{Setup}

Consider a cycle between $a_\mathrm{min}$ and $a_\mathrm{max}$
($0 < a_\mathrm{min} < a_\mathrm{max} < \infty$)
with ideal parallel plates at $T = 0$.

The potential (per unit area) is:
\begin{equation}
  V(a) = \frac{E(a)}{A} = -\frac{\pi^2}{720\, a^3}.
  \label{eq:potential}
\end{equation}

\subsection{Closing work ($a_\mathrm{max} \to a_\mathrm{min}$)}

The work done by the Casimir force during plate closing:
\begin{align}
  \frac{W_\mathrm{close}}{A}
  &= \int_{a_\mathrm{max}}^{a_\mathrm{min}} \frac{F(a)}{A}\, da
  = -\int_{a_\mathrm{max}}^{a_\mathrm{min}} \frac{dV}{da}\, da
  \nonumber \\
  &= -(V(a_\mathrm{min}) - V(a_\mathrm{max}))
  = V(a_\mathrm{max}) - V(a_\mathrm{min})
  \nonumber \\
  &= -\frac{\pi^2}{720}\left(\frac{1}{a_\mathrm{max}^3}
     - \frac{1}{a_\mathrm{min}^3}\right).
  \label{eq:wclose}
\end{align}

% SIGN CHECK: For a_min < a_max, 1/a_min^3 > 1/a_max^3,
% so (1/a_max^3 - 1/a_min^3) < 0,
% so W_close = -pi^2/720 * (negative) > 0.
% W_close > 0: the attractive force does positive work as plates close. CHECK.
%
% DIMENSION CHECK: [W/A] = [a^{-3}] = [length^{-3}]. CHECK.

\subsection{Opening work ($a_\mathrm{min} \to a_\mathrm{max}$)}

The work done by the Casimir force during plate opening:
\begin{align}
  \frac{W_\mathrm{open}}{A}
  &= \int_{a_\mathrm{min}}^{a_\mathrm{max}} \frac{F(a)}{A}\, da
  = V(a_\mathrm{min}) - V(a_\mathrm{max})
  \nonumber \\
  &= -\frac{\pi^2}{720}\left(\frac{1}{a_\mathrm{min}^3}
     - \frac{1}{a_\mathrm{max}^3}\right)
  = -\frac{W_\mathrm{close}}{A}.
  \label{eq:wopen}
\end{align}

% SIGN CHECK: W_open = -W_close < 0. The attractive force opposes
% opening, so the Casimir force does negative work (external agent
% must input energy). CHECK.

\subsection{Net cycle work}

\begin{equation}
  \boxed{
    \frac{W_\mathrm{net}}{A}
    = \frac{W_\mathrm{close}}{A} + \frac{W_\mathrm{open}}{A}
    = 0 \quad \text{(exact)}.
  }
  \label{eq:wnet}
\end{equation}

This is an exact algebraic cancellation, not a numerical coincidence.
It follows directly from $W_\mathrm{open} = -W_\mathrm{close}$,
which is a consequence of the force being conservative.

\subsection{Explicit numerical values}

For the canonical test cycle $a_\mathrm{min} = 100\,\mathrm{nm}$,
$a_\mathrm{max} = 1\,\mu\mathrm{m}$:

In natural units, the Casimir formulas $E/A = -\pi^2/(720\,a^3)$
and $F/A = -\pi^2/(240\,a^4)$ take $a$ in units of length.
Since we work with pure numbers (the code uses $a$ directly
in the formula), we can evaluate for any $a > 0$ with the
understanding that $a$ is measured in whatever length unit
we choose consistently.

Choosing $a$ in nanometers as the working unit:
\begin{align}
  a_\mathrm{min} &= 100, \quad a_\mathrm{max} = 1000 \quad
  \text{(in nm)} \\[6pt]
  \frac{W_\mathrm{close}}{A}
  &= -\frac{\pi^2}{720}\left(\frac{1}{1000^3} - \frac{1}{100^3}\right)
  \nonumber \\
  &= -\frac{\pi^2}{720}\left(10^{-9} - 10^{-6}\right)
  = -\frac{\pi^2}{720}\cdot(-999 \times 10^{-9})
  \nonumber \\
  &= +\frac{999\,\pi^2}{720 \times 10^9}
  \approx 1.369 \times 10^{-8} \quad [\text{nm}^{-3}].
  \label{eq:wclose-numerical}
\end{align}

% CROSS-CHECK: W_close > 0. CHECK.
% LIMITING CASE 1: a_min = a_max => W_close = 0. CHECK.
% LIMITING CASE 2: a_max -> infinity => W_close = pi^2/(720 a_min^3)
%   = -E(a_min)/A. This is all the stored vacuum energy being released. CHECK.

\subsection{Limiting cases}

\begin{enumerate}
  \item \textbf{Trivial cycle} ($a_\mathrm{min} = a_\mathrm{max}$):
    $W_\mathrm{close} = W_\mathrm{open} = W_\mathrm{net} = 0$. \checkmark

  \item \textbf{Infinite separation}
    ($a_\mathrm{max} \to \infty$):
    $W_\mathrm{close} \to -E(a_\mathrm{min})/A = \pi^2/(720\,a_\mathrm{min}^3)$.
    This is the total stored vacuum energy, confirming that bringing
    plates from infinity releases all the Casimir binding energy.
    $W_\mathrm{open} = -W_\mathrm{close}$: it costs exactly as much
    to separate the plates back to infinity. \checkmark

  \item \textbf{Small cycle}
    ($a_\mathrm{max} = a_\mathrm{min} + \delta$, $\delta \ll a_\mathrm{min}$):
    \begin{align}
      W_\mathrm{close} &\approx -F(a_\mathrm{min})\cdot\delta
      = \frac{\pi^2}{240\,a_\mathrm{min}^4}\,\delta > 0.
    \end{align}
    This is just ``force $\times$ displacement'' for an infinitesimal
    cycle, confirming consistency with the force definition. \checkmark
\end{enumerate}

%% ============================================================
\section{What Breaks Conservativeness}
%% ============================================================

The conservativeness of the Casimir force fails when any of
A1--A4 is violated:

\begin{enumerate}
  \item \textbf{Violation of A1} (material response changes during cycle):
    If $\varepsilon(\omega)$ changes between the closing and opening
    legs (e.g., via an external field or phase transition), then
    $F_\mathrm{close}(a) \neq F_\mathrm{open}(a)$, and the cycle
    integral need not vanish. This is the \emph{boundary condition
    change loophole} analyzed in Phase 05.

  \item \textbf{Violation of A2} (non-quasi-static, dynamic Casimir):
    If plates move at relativistic speeds, real photons are created
    (dynamic Casimir effect). The system is no longer in the
    equilibrium state at each $a$, and the equilibrium force law
    does not apply. Analyzed in Phase 06.

  \item \textbf{Violation of A3} (non-isothermal with $T$-dependent
    material response): If $T$ changes and $\varepsilon(\omega,T)$
    changes accordingly, the effective force law varies along the
    cycle. Proper treatment requires tracking entropy and heat flows.

  \item \textbf{Violation of A4} (topology change):
    Creating or destroying a cavity (e.g., merging two plates into
    one or introducing a third plate) changes the topology of the
    configuration space. The potential $V(a)$ on the old space is
    no longer the relevant quantity. Analyzed in Phase 05.
\end{enumerate}

%% ============================================================
\section{Energy Budget for a Closed $T=0$ Cycle}
%% ============================================================

For a closed isothermal cycle at $T = 0$ with conditions A1--A4:

\begin{center}
\begin{tabular}{l c c}
\hline
Quantity & Symbol & Value \\
\hline
Vacuum energy change & $\Delta E_\mathrm{vac}$ & $0$ (closed cycle) \\
Work by Casimir force (closing) & $W_\mathrm{close}$ & $>0$ \\
Work by Casimir force (opening) & $W_\mathrm{open}$ & $-W_\mathrm{close} < 0$ \\
Net work & $W_\mathrm{net}$ & $0$ (exact) \\
Heat exchanged & $Q$ & $0$ ($T = 0$) \\
\hline
Energy conservation & \multicolumn{2}{c}{$W_\mathrm{net} = \Delta E_\mathrm{vac} + Q = 0$} \\
\hline
\end{tabular}
\end{center}

The energy budget closes exactly. No net energy is extracted from
the vacuum over a complete cycle.

\end{document}
